\textbf{Анотація.} Цифрова трансформація наукового адміністрування закладів вищої освіти стримується фрагментацією наукометричних даних між ORCID, Scopus, Web of Science та OpenAlex. Неузгодженість ідентифікаторів, транслітерацій і метаданих призводить до 15--25\,\% дублікатів у типових звітах, а підготовка звітності потребує значної ручної верифікації. Мета статті --- обґрунтувати модель і механізми інтеграції GenAI та автономних агентів для проактивного управління науковою діяльністю ЗВО засобами ІКТ. Запропоновано ResearchIQ --- багатокористувацький вебзастосунок із React-інтерфейсом, Node.js API, tenant-сховищем даних і агентним шаром на основі \texttt{pi-agent-core}. Агент отримує локальний знімок даних ЗВО, викликає типізовані інструменти пошуку, ранжування, порівняння кафедр, побудови профілів, глобального пошуку праць і формування звітів, але не ухвалює критичних рішень щодо тотожності публікацій. Ядром обробки є «золотий запис» (Golden Record) --- конвеєр розв'язання сутностей із канонізацією метаданих, блокуванням, парним логістичним зіставленням, кластеризацією, провенансом і смугою перевірки людиною. Архітектуру спроєктовано за рамкою TRACE: детерміноване ядро, класичне машинне навчання та мовну модель-валідатор рознесено в окремі шари, а ощадливе використання LLM для структурованих даних виражено коефіцієнтом обчислювальної ощадливості. Модель апробовано на відтворюваному синтетичному бенчмарку, що імітує реалістичний український ЗВО (300 науково-педагогічних працівників, 20 кафедр, близько 5 факультетів) без персональних даних; статистичну стійкість оцінено методом Монте-Карло (K=30). Кількісну апробацію зосереджено на детермінованому ядрі дедуплікації; агентний шар подано як концептуальний природномовний інтерфейс, що не бере участі в обчисленні метрик. Golden Record досяг F1 76,8\,\% проти 59,7\,\% для DOI-only та 53,5\,\% для Title+Year. Виправлено дефект нормалізатора, що знищував 100\,\% україномовних назв публікацій. Аналітична модель оцінює ручну працю в 429,1 год/рік на ЗВО проти 2,11 с машинного обчислення. Результати позиціонують підхід як реалістичну стратегію цифрової трансформації наукового адміністрування.

\textbf{Ключові слова:} ІКТ в управлінні освітою; цифрова трансформація наукового адміністрування; генеративний штучний інтелект; автономний агент (MCP); золотий запис і розв'язання сутностей; проактивне управління науковою діяльністю ЗВО.

\bigskip

\textbf{Abstract.} The digital transformation of research administration in higher education institutions is hindered by fragmented scientometric data distributed across ORCID, Scopus, Web of Science, and OpenAlex. Inconsistent identifiers, transliterations, and metadata cause 15--25\,\% of records in typical reports to be duplicates and require substantial manual verification. The aim of the article is to substantiate a model and integration mechanisms for GenAI and autonomous agents that support proactive HEI research management through ICT. We propose ResearchIQ, a multi-tenant web application with a React interface, a Node.js API, tenant data storage, and an agent layer based on \texttt{pi-agent-core}. The agent receives a local HEI data snapshot and invokes typed tools for search, ranking, department comparison, faculty profiles, global work search, and report generation, but it does not make critical publication-identity decisions. The data-processing core is the Golden Record --- an Entity Resolution pipeline with metadata canonicalisation, blocking, pairwise logistic matching, clustering, provenance tracking, and a human-in-the-loop review band. The architecture follows the TRACE framework: a deterministic core, classical machine learning, and an LLM validator are separated into distinct layers, and parsimonious LLM use for structured-data tasks is expressed by the Computational Parsimony Ratio. The model is evaluated on a reproducible synthetic benchmark simulating a realistic Ukrainian HEI (300 academic staff, 20 departments, about 5 faculties) without personal data; statistical robustness is assessed by Monte Carlo simulation (K=30). The quantitative evaluation targets the deterministic deduplication core; the agent layer is a conceptual natural-language interface and is not itself benchmarked here. The Golden Record pipeline achieves an F1 score of 76.8\,\% versus 59.7\,\% for DOI-only and 53.5\,\% for Title+Year. A critical normaliser defect that destroyed 100\,\% of Ukrainian-language publication titles was corrected. An analytical efficiency model estimates manual effort at 429.1 hours per HEI-year compared with 2.11 seconds of machine computation. The findings position the approach as a realistic ICT strategy for transforming HEI research administration.

\textbf{Keywords:} ICT in education management; digital transformation of research administration; generative artificial intelligence; autonomous agent (MCP); Golden Record and Entity Resolution; proactive management of HEI research activity.
